\documentclass[portuguese,12pt,oneside,a4paper]{book}

% ---------------------------------------------------------------------------- %
% Pacotes 
\usepackage[T1]{fontenc}
\usepackage[brazil]{babel}              % define a linguagem padrão
\usepackage[utf8x]{inputenc}

% --- AJUSTES ABNT DE FONTE E PARÁGRAFO ---
\usepackage{mathptmx}                   % Usa a fonte Times New Roman (Padrão ABNT)
\setlength{\parindent}{1.25cm}          % Recuo da primeira linha do parágrafo (ABNT)

\usepackage[pdftex]{graphicx}           % usamos arquivos pdf/png como figuras
\usepackage{setspace}                   % espaçamento flexível
\usepackage{makeidx}                    % índice remissivo
\usepackage[nottoc]{tocbibind}          % acrescentamos a bibliografia no Table of Contents
\usepackage{courier}                    
\usepackage{type1cm}                    
\usepackage{listings}                   
\usepackage{titletoc}

% --- AJUSTE ABNT PARA OS TÍTULOS DOS CAPÍTULOS ---
\usepackage{titlesec}
\titleformat{\chapter}[hang]{\normalfont\bfseries\normalsize}{\thechapter}{1em}{\MakeUppercase}
\titlespacing*{\chapter}{0pt}{-30pt}{1.5cm}
\titleformat{\section}[hang]{\normalfont\bfseries\normalsize}{\thesection}{1em}{}
\titleformat{\subsection}[hang]{\normalfont\normalsize}{\thesubsection}{1em}{}

\usepackage[fixlanguage]{babelbib}
\usepackage[font=small,format=plain,labelfont=bf,up,textfont=it,up]{caption}
\usepackage[usenames,svgnames,dvipsnames]{xcolor}

% --- MARGENS ABNT CORRIGIDAS ---
% Superior: 3cm, Esquerda: 3cm, Inferior: 2cm, Direita: 2cm
\usepackage[a4paper,top=3.0cm,bottom=2.0cm,left=3.0cm,right=2.0cm]{geometry} 

\usepackage[pdftex,plainpages=false,pdfpagelabels,pagebackref,colorlinks=true,citecolor=black,linkcolor=black,urlcolor=black,filecolor=black,bookmarksopen=true]{hyperref} 
\usepackage[all]{hypcap}                

\fontsize{60}{62}\usefont{OT1}{cmr}{m}{n}{\selectfont}
\usepackage{color} 
\usepackage{xcolor}
\usepackage{amsmath} 
\usepackage{pdflscape} 
\usepackage{longtable} 
\usepackage{multicol}
\usepackage{enumitem} 
\usepackage{pdfpages}

% ---------------------------------------------------------------------------- %
% Cabeçalhos
\usepackage{fancyhdr}
\pagestyle{fancy}
\fancyhf{}
\renewcommand{\chaptermark}[1]{\markboth{\MakeUppercase{#1}}{}}
\renewcommand{\sectionmark}[1]{\markright{\MakeUppercase{#1}}{}}
\renewcommand{\headrulewidth}{0pt}

% ---------------------------------------------------------------------------- %
\graphicspath{{./figuras/}}             % O LaTeX vai procurar as imagens na pasta "figuras"
\frenchspacing                          
\urlstyle{same}                         
\makeindex                              
\raggedbottom                           
\normalsize

% ---------------------------------------------------------------------------- %
% Opções de listing usados para o código fonte
\lstset{ 
language=Java,                  
basicstyle=\footnotesize,       
numbers=left,                   
numberstyle=\footnotesize,      
stepnumber=1,                   
numbersep=5pt,                  
showspaces=false,               
showstringspaces=false,         
showtabs=false,                 
frame=single,                   
framerule=0.6pt,
tabsize=2,                      
captionpos=b,                   
breaklines=true,                
breakatwhitespace=false,        
escapeinside={\%*}{*)},         
backgroundcolor=\color[rgb]{1.0,1.0,1.0}, 
rulecolor=\color[rgb]{0.8,0.8,0.8},
extendedchars=true,
xleftmargin=10pt,
xrightmargin=10pt,
framexleftmargin=10pt,
framexrightmargin=10pt
}

% ABNT CITAÇÕES %%%%%%%%%%%%%%%%%%%%%%%%%%%%%%%%%%%%%%%%%%%%%%%%%%%%%%%%%%%%%%%%%%%%
\usepackage[alf,
abnt-etal-text=emph,
abnt-and-type=e,
bibjustif,
abnt-etal-cite=3]{abntex2cite} 
\usepackage{cite}

\makeatletter
\renewcommand\@biblabel[1]{[#1]~}
\renewcommand\citepunct{;~}
\makeatother
\citebrackets[] 
%%%%%%%%%%%%%%%%%%%%%%%%%%%%%%%%%%%%%%%%%%%%%%%%%%%%%%%%%%%%%%%%%%%%%%%%%%%%%%%%%%%%

% ---------------------------------------------------------------------------- %
% Corpo do texto
\begin{document}
\frontmatter 
\fancyhead[RO]{{\footnotesize\rightmark}\hspace{2em}\thepage}
\setcounter{tocdepth}{2}
\fancyhead[LE]{\thepage\hspace{2em}\footnotesize{\leftmark}}
\fancyhead[RE,LO]{}
\fancyhead[RO]{{\footnotesize\rightmark}\hspace{2em}\thepage}

\onehalfspacing  % espaçamento 1.5 obrigatório na ABNT

% ---------------------------------------------------------------------------- %
% Capa
\thispagestyle{empty}
\begin{center}
\includegraphics[scale=0.4]{ifpb}\\ 
\Large Instituto Federal de Educação, Ciência e Tecnologia da Paraíba\\
Campus Campina Grande\\
Coordenação do Curso Superior de Tecnologia em Telemática\\

\vspace*{4cm}

    \textbf{\LARGE{DESENVOLVIMENTO DE UM SISTEMA MULTIPLATAFORMA PARA ANÁLISE E INTERPRETAÇÃO AUTOMATIZADA DE ELETROCARDIOGRAMAS (ECG) APLICADO À PRÁTICA CLÍNICA}}\\
    
\vspace*{2cm}

\Large ACHILLES ROBSON VASCONCELOS

\vskip 4cm

Orientador: Nome do Orientador\\

\vskip 1.85cm

\normalsize Campina Grande, PB, 2026 \\

\end{center}


\newpage

% Folha de Rosto
\thispagestyle{empty}

\begin{center}
\includegraphics[scale=0.4]{ifpb}\\ 
    \Large Instituto Federal de Educação, Ciência e Tecnologia da Paraíba\\
    Campus Campina Grande\\
    Coordenação do Curso Superior de Tecnologia em Telemática\\

    \vspace*{3cm}

    \textbf{\LARGE{DESENVOLVIMENTO DE UM SISTEMA MULTIPLATAFORMA PARA ANÁLISE E INTERPRETAÇÃO AUTOMATIZADA DE ELETROCARDIOGRAMAS (ECG) APLICADO À PRÁTICA CLÍNICA}}\\
    
    \vspace*{2cm}
    \Large{ACHILLES ROBSON VASCONCELOS}
    
    \vskip 1.5cm
\end{center}   
    
    \begin{flushright} 
        \begin{minipage}[left]{0.5\linewidth}
            Monografia apresentada à Coordenação do Curso de Telemática do IFPB, como requisito parcial para conclusão do curso de Tecnologia em Telemática.                       
        \end{minipage}
    \end{flushright}                               

\begin{center}  
    \vskip 2cm
    Orientador: Nome do Orientador\\
    \vskip 1.5cm
    \normalsize{Campina Grande, PB, 2026}
\end{center}

\newpage

% ---------------------------------------------------------------------------- %
% Ficha de Aprovação
\newpage
\thispagestyle{empty}
    \begin{center}
        
        \textbf{\LARGE{DESENVOLVIMENTO DE UM SISTEMA MULTIPLATAFORMA PARA ANÁLISE E INTERPRETAÇÃO AUTOMATIZADA DE ELETROCARDIOGRAMAS (ECG) APLICADO À PRÁTICA CLÍNICA}}\\

        \vspace*{3cm}

        \textbf{ACHILLES ROBSON VASCONCELOS}
    
        \vskip 2cm

        \rule{10cm}{0.1mm}\\
        Nome do Orientador\\
        Orientador
        
        \vskip 2cm

        \rule{10cm}{0.1mm}\\
        Mikael Marinho Oliveira\\
        Coorientador
        
        \vskip 2cm
                
        \rule{10cm}{0.1mm}\\
        Nome do Membro da Banca 1\\
        Membro da Banca
            
        \vskip 2cm
        
        \rule{10cm}{0.1mm}\\
        Nome do Membro da Banca 2\\
        Membro da Banca
            
        \vskip 1cm
            
        Campina Grande, Paraíba, Brasil\\
        2026

    \end{center}
\pagebreak

\newpage
\thispagestyle{empty}
        \vspace*{20.0cm}
        \begin{flushright}
            À minha esposa e ao meu filho, fontes inesgotáveis de apoio e inspiração.
        \end{flushright}

\pagebreak

% ---------------------------------------------------------------------------- %
% Agradecimentos
\chapter*{Agradecimentos}

Gostaria de expressar minha profunda gratidão a todos que colaboraram direta e indiretamente para a realização deste trabalho. 

Aos professores do Instituto Federal da Paraíba (IFPB), pelos ensinamentos e pela sólida base na área de Telemática. 

Aos colegas e médicos das unidades de pronto atendimento, cujas rotinas e desafios práticos serviram de inspiração para que eu pudesse unir meus conhecimentos de Técnico de Enfermagem com o desenvolvimento de software.

\pagebreak


% ---------------------------------------------------------------------------- %
% Resumo
\chapter*{Resumo}
As doenças cardiovasculares representam a principal causa de morbimortalidade global, tornando o Eletrocardiograma (ECG) um exame crítico nas emergências. Contudo, a leitura precisa dos traçados apresenta desafios para muitos profissionais da linha de frente, gerando atrasos em condutas tempo-dependentes. Este trabalho apresenta o \textit{ECG Analyzer}, um Sistema de Apoio à Decisão Clínica multiplataforma desenvolvido para mitigar essa lacuna assistencial. O sistema foi construído utilizando React.js no \textit{front-end} e Node.js no \textit{back-end}, operando diretamente no navegador sem necessidade de instalações. A solução destaca-se pela capacidade de ler exames impressos através de fotografias utilizando a \textit{Canvas API}, e processar os dados por meio de um Motor de Diagnóstico proprietário. O algoritmo alcançou uma acurácia média de 80\% na análise de traçados e é capaz de cruzar dados do ECG com mais de 14 sinais vitais e laboratoriais, rastreando mais de 100 diagnósticos e 30 síndromes multissistêmicas. Conclui-se que a aplicação da Telemática aliada à prática da enfermagem e medicina resulta em uma ferramenta viável, capaz de gerar laudos em PDF e democratizar o acesso a triagens cardiológicas de alta qualidade em ambientes hospitalares.

\noindent \textbf{Palavras-chave:} eletrocardiograma; inteligência artificial; telemática; Sistemas de Apoio à Decisão Clínica; React.js.

\pagebreak

% ---------------------------------------------------------------------------- %
% Abstract
\chapter*{Abstract}
Cardiovascular diseases represent the main cause of global morbidity and mortality, making the Electrocardiogram (ECG) a critical exam in emergencies. However, the precise reading of the tracings presents challenges for many frontline professionals, causing delays in time-dependent interventions. This work presents the \textit{ECG Analyzer}, a cross-platform Clinical Decision Support System developed to mitigate this care gap. The system was built using React.js on the \textit{front-end} and Node.js on the \textit{back-end}, operating directly in the browser without the need for installations. The solution stands out for its ability to read printed exams through photographs using the \textit{Canvas API}, and process the data through a proprietary Diagnostic Engine. The algorithm achieved an average accuracy of 80\% in tracing analysis and is capable of cross-referencing ECG data with more than 14 vital and laboratory signs, tracking more than 100 diagnoses and 30 multisystemic syndromes. It is concluded that the application of Telematics combined with nursing and medical practice results in a viable tool, capable of generating PDF reports and democratizing access to high-quality cardiological screening in hospital environments.

\noindent \textbf{Keywords:} electrocardiogram; artificial intelligence; telematics; Clinical Decision Support Systems; React.js.

% ---------------------------------------------------------------------------- %
% Sumário
\tableofcontents    % imprime o sumário

% ---------------------------------------------------------------------------- %
\chapter*{Lista de Siglas}
\begin{tabular}{ll}
    API     & \textit{Application Programming Interface} (Interface de Programação de Aplicações) \\
    CSV     & \textit{Comma-Separated Values} (Valores Separados por Vírgula) \\
    DCV     & Doenças Cardiovasculares \\
    ECG     & Eletrocardiograma \\
    IAM     & Infarto Agudo do Miocárdio \\
    JSON    & \textit{JavaScript Object Notation} (Notação de Objetos JavaScript) \\
    PDS     & Processamento Digital de Sinais \\
    SADC    & Sistema de Apoio à Decisão Clínica \\
    UI      & \textit{User Interface} (Interface de Usuário) \\
    UX      & \textit{User Experience} (Experiência do Usuário) \\
    VD / VE & Ventrículo Direito / Ventrículo Esquerdo \\
\end{tabular}

% ---------------------------------------------------------------------------- %
% Listas de figuras e tabelas criadas automaticamente
\listoffigures            
\listoftables            

% ---------------------------------------------------------------------------- %
% Capítulos do trabalho
\mainmatter

% cabeçalho para as páginas de todos os capítulos
\fancyhead[RE,LO]{\thesection}
\onehalfspacing            % espaçamento um e meio

% ============================================================================ %
% CAPÍTULO 1
% ============================================================================ %
\chapter{INTRODUÇÃO}

As doenças cardiovasculares (DCV) representam, de acordo com a Organização Mundial da Saúde (OMS), a principal causa de morbimortalidade em escala global, gerando impactos significativos tanto na qualidade de vida dos pacientes quanto nos custos operacionais dos sistemas de saúde públicos e privados. Diante da natureza aguda e, muitas vezes, fatal de eventos como o Infarto Agudo do Miocárdio (IAM) e as arritmias malignas, o diagnóstico precoce e preciso constitui um fator determinante para o prognóstico e a sobrevida do paciente. 

Nesse cenário de urgência e emergência, o Eletrocardiograma (ECG) consolida-se como o exame padrão-ouro para a triagem e o diagnóstico inicial. Trata-se de um método não invasivo, de baixo custo e amplamente disponível nas unidades de saúde. No entanto, a eficácia diagnóstica do ECG não reside apenas na captação do sinal elétrico, mas depende intrinsecamente da capacidade e da experiência do profissional de saúde na interpretação dos traçados. A análise de um ECG exige o reconhecimento de padrões morfológicos complexos e variações sutis nos intervalos das ondas, uma habilidade que exige treinamento especializado e contínuo.

A observação empírica direta na prática clínica diária – especificamente na rotina da enfermagem e no atendimento médico de urgência de linha de frente – revela um gargalo estrutural significativo. Constata-se que muitos médicos plantonistas não especialistas (como clínicos gerais emergencistas) e profissionais das equipes de triagem apresentam dificuldades técnicas na interpretação autônoma e precisa de eletrocardiogramas complexos. Essa insegurança ou falta de expertise na leitura dos sinais pode resultar em atrasos no início de terapias trombolíticas, encaminhamentos incorretos ou falhas na identificação de alterações isquêmicas sutis, comprometendo condutas clínicas que são tempo-dependentes.

Motivado por essa lacuna assistencial, e compreendendo que a Telemática aliada à Inteligência Artificial possui o potencial de transformar dados brutos em informações clinicamente acionáveis, este trabalho propõe uma intervenção tecnológica. A convergência entre a ciência da computação e a prática assistencial permite a criação de Sistemas de Apoio à Decisão Clínica (SADC) capazes de mitigar o erro humano e otimizar o fluxo de triagem hospitalar. 

Dessa forma, a presente pesquisa tem como foco explorar o processamento digital de sinais biomédicos e o desenvolvimento de software para entregar uma solução acessível que atue como uma ferramenta de segunda opinião, conferindo maior segurança à equipe multidisciplinar no momento do diagnóstico.

\section{Objetivos}

\subsection{Objetivo Geral}
O objetivo geral deste trabalho é desenvolver um sistema computacional multiplataforma (acessível via dispositivos móveis e computadores) capaz de realizar a leitura, o processamento e a análise automatizada de traçados de Eletrocardiograma (ECG), fornecendo suporte à decisão clínica com uma acurácia média de 80\%.

\subsection{Objetivos Específicos}
Para alcançar o objetivo geral, foram traçados os seguintes objetivos específicos:
\begin{itemize}
    \item Desenvolver e otimizar um Motor de Diagnóstico em Node.js voltado para a filtragem de sinais e identificação de padrões cardíacos baseados em diretrizes médicas.
    \item Implementar uma arquitetura de software robusta, utilizando \textit{back-end} (Express) para o processamento de dados e \textit{front-end} em JavaScript (React) para garantir responsividade e usabilidade.
    \item Integrar rotinas computacionais capazes de cruzar achados do ECG com parâmetros de sinais vitais e laboratoriais para diagnóstico de síndromes multissistêmicas.
    \item Avaliar e validar o desempenho do algoritmo, assegurando que o sistema possa ser utilizado de forma fluida tanto em estações de trabalho quanto à beira-leito através de dispositivos móveis.
\end{itemize}

\section{Justificativa}
A relevância deste projeto justifica-se sob as óticas social, científica e tecnológica. Socialmente, o sistema propõe-se a democratizar o acesso a análises cardiológicas de qualidade, especialmente em Unidades de Pronto Atendimento (UPAs) e hospitais de pequeno porte onde a figura do médico cardiologista nem sempre está presente no plantão. 

Do ponto de vista científico e tecnológico, o trabalho demonstra a viabilidade da criação de algoritmos de processamento de exames integrados a marcadores biológicos (como eletrólitos e troponina). A transformação de um modelo teórico em um aplicativo prático e multiplataforma reforça o papel da Telemática como ponte entre o desenvolvimento de software e a inovação médica.

% ============================================================================ %
% CAPÍTULO 2
% ============================================================================ %
\chapter{FUNDAMENTAÇÃO TEÓRICA}

\section{O Eletrocardiograma na Prática Clínica}
O ECG registra a atividade elétrica do coração, refletindo eventos de despolarização e repolarização celular, representados classicamente pelas ondas P, complexo QRS e onda T. A identificação de alterações morfológicas ou de ritmo nesses componentes é fundamental para detectar anomalias, desde bloqueios atrioventriculares até síndromes coronarianas agudas \cite{hannun2019}.

\section{Sistemas de Apoio à Decisão Clínica e Telemática}
A aplicação de algoritmos computacionais e de sistemas baseados em regras ou inteligência artificial no processamento de biossinais tem crescido exponencialmente. Sistemas de Apoio à Decisão Clínica (SADC) processam grandes volumes de variáveis clínicas para identificar padrões imperceptíveis ou auxiliar a triagem rápida. Na Telemática, a transmissão e o processamento de dados médicos exigem confiabilidade e tempo de resposta baixo \cite{hannun2019}.

\section{Ecossistema JavaScript em Aplicações Médicas}
Historicamente, o Processamento Digital de Sinais (PDS) dependia de linguagens compiladas pesadas. Atualmente, o ecossistema JavaScript provou sua maturidade. A utilização do React.js permite a construção de interfaces de usuário (UI) reativas em múltiplas plataformas, enquanto o Node.js no \textit{back-end} fornece processamento assíncrono altamente eficiente, ideal para calcular matrizes de dados extraídas de exames médicos sem travar o dispositivo do usuário final \cite{react2026, nodejs2026}.

% ============================================================================ %
% CAPÍTULO 3
% ============================================================================ %
\chapter{METODOLOGIA}

O projeto foi caracterizado como uma pesquisa aplicada com desenvolvimento experimental, conduzida ao longo de nove meses. O objetivo metodológico central foi construir um sistema computacional altamente responsivo e clinicamente útil, capaz de processar não apenas os dados brutos de um ECG, mas também correlacioná-los com sinais vitais e marcadores laboratoriais. 

Para estruturar o desenvolvimento, mitigar problemas técnicos e garantir entregas funcionais, a construção do sistema foi dividida em cinco etapas iterativas e incrementais:

\section{Primeira Etapa: Prototipação e Estruturação de Dados Iniciais}
O projeto teve início com a modelagem da entrada de dados. O objetivo inicial era validar a leitura e a extração de métricas básicas de traçados de ECG. 
Foi enfrentada uma dificuldade inerente em obter sinais contínuos reais no início do desenvolvimento para testar os algoritmos. Para solucionar isso, foi implementado um módulo de \textit{upload} para arquivos padronizados (CSV e JSON) e criado um gerador de dados simulados. Isso permitiu à ferramenta plotar gráficos preliminares contendo taxa de amostragem e milivoltagens.

\section{Segunda Etapa: Processamento de Imagens e Extração de Sinais}
Na transição para o uso prático, identificou-se que médicos e enfermeiros geralmente não possuem o arquivo digital do ECG à beira-leito, mas sim o papel impresso, do qual tiram fotos.
A solução encontrada foi desenvolver um método de extração de dados utilizando a \textit{Canvas API} nativa do navegador. Essa abordagem permitiu ler o contraste das imagens digitais do traçado (fundo claro com traço escuro), traduzindo as variações de pixels em matrizes de dados numéricos que puderam ser plotados interativamente através da biblioteca Chart.js.

\section{Terceira Etapa: Desenvolvimento da Arquitetura Web e Interface}
O \textit{front-end} foi consolidado utilizando o \textit{framework} React.js 18 em conjunto com o \textit{bundler} Vite e a biblioteca de estilização Tailwind CSS 3. Essa escolha arquitetônica garantiu o desenvolvimento de uma interface responsiva, com suporte nativo a \textit{Dark Mode} (essencial para ambientes hospitalares noturnos) e formulários dinâmicos para entrada de variáveis clínicas.

\section{Quarta Etapa: Motor de Diagnóstico e \textit{Back-end}}
O escopo do projeto foi ampliado para um Sistema de Apoio à Decisão Clínica completo. A solução foi estruturar o \textit{back-end} em Node.js com Express, onde foi encapsulado um "Motor de Diagnóstico" proprietário. Este motor foi programado para cruzar as características morfológicas do ECG com 14 parâmetros de sinais vitais, resultando em uma árvore de decisão capaz de classificar mais de 100 diagnósticos distintos e identificar mais de 30 síndromes multissistêmicas.

\section{Quinta Etapa: Geração de Laudos e Histórico}
Foi implementado um sistema de persistência de dados no navegador (\textit{localStorage}), garantindo que o histórico das análises ficasse salvo de forma segura. Adicionalmente, desenvolveu-se um gerador de PDF profissional embutido no React, capaz de compilar os diagnósticos priorizados por severidade, a interpretação clínica e um espaço para a assinatura do profissional.

% ============================================================================ %
% CAPÍTULO 4
% ============================================================================ %
\chapter{RESULTADOS}

O principal resultado deste trabalho culminou na entrega do \textit{ECG Analyzer}\footnote{O aplicativo encontra-se disponível para uso \textit{online}. O código-fonte aberto pode ser acessado no repositório GitHub do autor: \url{https://github.com/Achillesdev/ECG-Analyzer-react}.}, uma aplicação web funcional, escalável e de alta responsividade. A evolução do sistema permitiu consolidar uma arquitetura moderna e de alto desempenho baseada em React.js e Node.js.

O marco central da pesquisa foi o sucesso na implementação do Motor de Diagnóstico. Os testes revelaram que o algoritmo atinge uma acurácia média de 80\% na detecção de anomalias puras nos traçados de ECG. O sistema superou o escopo técnico inicial ao processar, de forma cruzada, os dados morfológicos do ECG com 14 parâmetros distintos de sinais vitais e laboratoriais (como saturação de oxigênio, níveis de potássio e marcadores de troponina). Como resultado, a plataforma classifica mais de 100 diagnósticos isolados e aponta mais de 30 síndromes multissistêmicas.

Na usabilidade hospitalar, a integração da \textit{Canvas API} permitiu que o sistema extraísse dados a partir de fotografias comuns de exames (PNG e JPG), eliminando a necessidade de arquivos brutos. O sistema roda de forma fluida tanto em computadores quanto em \textit{smartphones} à beira-leito. Destacam-se as seguintes entregas funcionais:
\begin{itemize}
    \item \textbf{Geração de Laudos Profissionais:} Emissão de PDF estruturado com nível de risco visual, resumo clínico e campo de assinatura com CRM.
    \item \textbf{Gestão de Histórico:} Armazenamento automático via \textit{localStorage}, permitindo acesso rápido a exames prévios.
    \item \textbf{Acessibilidade Visual:} Interface reativa com \textit{Dark Mode} adaptativo.
\end{itemize}

% ============================================================================ %
% CAPÍTULO 5
% ============================================================================ %
\chapter{DISCUSSÃO}

A acurácia de 80\% alcançada pelo sistema está em consonância com estudos iniciais de algoritmos de triagem e ferramentas comerciais (como \textit{wearables} de última geração), o que demonstra a validade científica da abordagem algorítmica e baseada em regras clínicas estruturadas pelo desenvolvedor. 

Observa-se que estudos com Redes Neurais Profundas atingem margens superiores a 90\%, porém exigem alto poder computacional e conectividade avançada. O modelo proposto (\textit{ECG Analyzer}) destaca-se justamente pela leveza de execução no navegador e pelo cruzamento inteligente com os dados laboratoriais. A integração de variáveis como "Troponina" e "Potássio" na equação diagnóstica replica o raciocínio holístico humano, prevenindo interpretações isoladas do gráfico de ECG que poderiam gerar falsos positivos.

% ============================================================================ %
% CAPÍTULO 6
% ============================================================================ %
\chapter{CONCLUSÃO}

O presente projeto demonstra, de forma prática e conclusiva, a viabilidade e a extrema relevância da aplicação da Telemática e do desenvolvimento web no suporte ao diagnóstico médico. A conversão de uma percepção clínica empírica – a dificuldade na leitura de ECGs – resultou no \textit{ECG Analyzer}.

A implementação do Motor de Diagnóstico proprietário superou as expectativas ao correlacionar dados morfológicos extraídos de imagens com 14 parâmetros de sinais vitais, rastreando mais de 100 diagnósticos e 30 síndromes complexas. 

O sistema possui grande potencial para ser implantado em Unidades de Pronto Atendimento e clínicas do interior, democratizando o acesso a análises cardiológicas especializadas. Por fim, este trabalho evidencia o poder da interdisciplinaridade: ao unir a vivência assistencial técnica da enfermagem com a engenharia de software \textit{Full Stack}, foi possível entregar uma ferramenta que otimiza o tempo de resposta em cenários de emergência, colaborando para a preservação da vida.

% ---------------------------------------------------------------------------- %
% Apêndices e/ou Anexos
\appendix
\renewcommand{\appendixname}{Apêndice} 
\setcounter{chapter}{0} 

\chapter{Telas e Acesso ao Sistema}

Neste apêndice, documenta-se a interface gráfica final do \textit{ECG Analyzer}, demonstrando sua responsividade, usabilidade clínica e a geração de laudos. O acesso ao sistema em ambiente de produção pode ser realizado através do repositório público fornecido no Capítulo de Resultados.

A interface foi projetada com foco na experiência do usuário (UX) em ambientes hospitalares, adotando o \textit{Dark Mode} por padrão para reduzir a fadiga visual durante plantões noturnos.

% --- Figura 1: Tela Principal / Dashboard ---
\begin{figure}[htbp]
    \centering
    \caption{Tela inicial e painel de \textit{upload} de exames (Dashboard)}
    \includegraphics[width=0.9\textwidth]{tela-principal.jpeg}
    \caption*{Fonte: Elaborado pelo autor (2026).}
    \label{fig:tela-principal}
\end{figure}

A Figura \ref{fig:tela-principal} apresenta a tela inicial do sistema, onde o profissional de saúde pode realizar o \textit{upload} da imagem do ECG impresso, bem como inserir os dados vitais e laboratoriais do paciente para análise cruzada.

\vspace{1cm} 

% --- Figura 2: Tela de Sinais Vitais e Diagnóstico ---
\begin{figure}[htbp]
    \centering
    \caption{Interface de inserção de parâmetros clínicos e resultados da análise}
    \includegraphics[width=0.9\textwidth]{tela-diagnostico.jpeg}
    \caption*{Fonte: Elaborado pelo autor (2026).}
    \label{fig:tela-diagnostico}
\end{figure}

Na Figura \ref{fig:tela-diagnostico}, observa-se a capacidade multissistêmica da ferramenta. O motor de diagnóstico exibe a classificação de severidade (através de \textit{badges} visuais) e lista as síndromes e anomalias detectadas pelo cruzamento do traçado com exames de sangue e eletrólitos.

\vspace{1cm}

% --- Figura 3: Versão Mobile ---
\begin{figure}[htbp]
    \centering
    \caption{Responsividade do sistema em dispositivos móveis (Mobile)}
    \includegraphics[width=0.4\textwidth]{tela-celular.png}
    \caption*{Fonte: Elaborado pelo autor (2026).}
    \label{fig:tela-celular}
\end{figure}

Garantindo a aplicabilidade à beira-leito, a Figura \ref{fig:tela-celular} demonstra a adaptação automática do \textit{layout} para telas menores, permitindo que a equipe de enfermagem ou médica utilize o sistema diretamente em \textit{smartphones} durante a triagem.

\vspace{1cm}

% --- Figura 4: Laudo em PDF ---
\begin{figure}[htbp]
    \centering
    \caption{Laudo profissional gerado automaticamente pelo sistema (PDF)}
    \includegraphics[width=0.8\textwidth]{tela-laudo.png}
    \caption*{Fonte: Elaborado pelo autor (2026).}
    \label{fig:tela-laudo}
\end{figure}

% ---------------------------------------------------------------------------- %
% Bibliografia
\backmatter \singlespacing   % espaçamento simples
\bibliographystyle{abntex2-alf}% citação bibliográfica
\nocite{achilles2026, hannun2019, react2026, nodejs2026}
\bibliography{bibliografia}  % associado ao arquivo: 'bibliografia.bib'

\end{document}